\section{Conclusion}

In this manuscript, we introduce the concept of 3D Egocentric Foundation Models that integrate egocentric sensor data for 3D scene understanding. We identify two core tasks for 3D EFMs---3D object detection and surface regression---and create a benchmarks for each task using high quality annotations of datasets captured with Project Aria glasses~\cite{engel2023project}. When evaluating these tasks over an entire sequence of egocentric data (as opposed to a single frame), existing methods exhibit poor 3D consistency in their predictions that leads to poor performance on the EFM3D benchmark. To address this, we design a simple but effective 3D backbone for egocentric data, \EFMmodel{}, that leverages semi-dense points and image features to produce a 3D voxel grid of features. \EFMmodel{} outperforms all other methods when evaluated on the proposed EFM3D benchmark.
%We train \EFMmodel{} using high quality simulated egocentric data from ASE~\cite{ase2024web}. Despite being trained on simulated data, our model significantly outperforms all other methods when evaluated on real-world egocentric data for both EFM3D tasks. 
The simplicity of this architecture underscores the effectiveness of the 3D inductive biases in \EVL{}. 
%We leave further architecture exploration and optimization to future work. 
We encourage the development of more sophisticated models that can exploit the richness of egocentric 3D data even more effectively, including the incorporation of dynamic scene understanding and user interaction modeling. Such modeling advancements could improve the performance even further and extend the applicability of 3D EFMs to a wider range of real-world scenarios.
